% !TeX root = ../../Book4.tex
% 纲要标题：### 6.5 四点标记球面的参数与退化
% 纲要位置：几何学计划纲要.md，第 3726 行。
% 正文状态：初稿；尚未正式审核。

% 纲要内容开始
% 
% 这里完成解析参数和三类退化的识别；跨任意代数基底的通用性证明在 V.34.6 进行。交比参数化的是模空间，不把它等同于其带标记的 Teichmüller 覆盖。
% 
% * 有序四点、Möbius 归一化与交比参数
% * $\mathbb P^1\setminus\{0,1,\infty\}$ 上的光滑四点族
% * 标记碰撞、长颈与结点极限
% * 两两分配标记的三种稳定边界类型
% * 有序标记与置换商、自同构的区别
% * 与 V.34.6 的代数通用族及完整紧化构造比较
% 
% ---
% 
% 纲要内容结束

\section{四点标记球面的参数与退化}
\label{g4:sec:0605}

球面本身没有复模参数，但在球面上指定四个有序的点后，第四个点相对于前三点的位置不能再由坐标变化消去。这个一维问题可以直接算完：光滑对象由一个交比描述，参数逼近三个缺失值时，碰撞的点分别落到新出现的球面分支上。下面同时保留参数、标记和局部方程，以便看清模空间中的边界点究竟表示什么。

\subsection{有序四点、Möbius 归一化与交比参数}
\label{g4:subsec:0605:01}
设 \(p_1,p_2,p_3,p_4\) 是 \(\mathbb P^1\) 上两两不同的点。这里的有序标记是四个点的编号；本章前面的一张拓扑曲面到复曲面的标记则包含道路和同位信息。四点的编号本身不记录这些道路。

\begin{proposition}
	\label{g4:prop:0605:cross-ratio}
	球面的双全纯自同构恰为 \(z\mapsto(az+b)/(cz+d)\)、\(ad-bc\ne0\)。任意三个不同的有序点有唯一的这类变换 \(N\)，满足 \(N(p_1)=0\)、\(N(p_2)=1\)、\(N(p_3)=\infty\)。数值
	\[
		\lambda=N(p_4)=\frac{(p_4-p_1)(p_2-p_3)}{(p_4-p_3)(p_2-p_1)}
	\]
	按无穷远处的极限解释，是有序四点球面在保持编号的双全纯同构下的完全不变量，其取值集为 \(B=\mathbb P^1\setminus\{0,1,\infty\}\)。
\end{proposition}
\begin{proof}
	先解释自同构的形式。设 \(F\) 为球面自同构，把唯一的 \(F^{-1}(\infty)\) 用一个分式线性坐标移到无穷远。此后 \(F\) 在有限平面全纯，并因无穷远处局部可逆而只有一阶极点。因此 \(|F(z)|\le C(1+|z|)\)，Cauchy 估计使二阶导数恒为零，\(F\) 是斜率非零的仿射函数。换回坐标得到所述分式形式。反之矩阵行列式不为零时，逆矩阵给出全纯逆。

	当前三点有限时，可取 \(N(z)=(z-p_1)(p_2-p_3)/((z-p_3)(p_2-p_1))\)；有无穷远点时取极限。若两变换均满足归一化，其复合逆固定 \(0,1,\infty\)：固定无穷远迫使它仿射，固定零和一再使它恒等。故归一化唯一。第四点不同于前三点等价于 \(\lambda\ne0,1,\infty\)。两组四点交比相同，当且仅当各自归一化后都是 \((0,1,\infty,\lambda)\)，因而可由保持编号的双全纯映射联系。
\end{proof}

这类分式变换称为 Möbius 变换。若将四个点看成球面配置空间的点，上述归一化在各坐标片上是全纯的，且分母只在被排除的碰撞处为零。因此交比不仅给出集合分类，也给出这个配置空间商的一维复坐标。我们记有序四点球面的模空间为 \(\mathcal M_{0,4}\)，便有 \(\mathcal M_{0,4}\simeq B\)。

\subsection{\texorpdfstring{\(\mathbb P^1\setminus\{0,1,\infty\}\) 上的光滑四点族}{P¹∖\{0,1,∞\} 上的光滑四点族}}
\label{g4:subsec:0605:02}
在交比坐标下，所有光滑对象同时出现在一个简单的族中：
\[
	\pi:\mathcal C=\mathbb P^1\times B\longrightarrow B,
	\qquad s_1(\lambda)=0,\quad s_2(\lambda)=1,\quad
	s_3(\lambda)=\infty,\quad s_4(\lambda)=\lambda.
\]
这里 \(s_i\) 是投影的截面，简写中省去了共同的第二坐标 \(\lambda\)。投影是固有的全纯下沉，四个截面处处不交。给定 \(\lambda\) 的纤维就是前一命题的标准代表。

\begin{proposition}
	\label{g4:prop:0605:family-normalization}
	设 \(T\) 为复流形，\(C\to T\) 是局部全纯平凡的球面族，带四个两两不交的有序全纯截面。则存在唯一全纯映射 \(\lambda:T\to B\)，以及唯一保持四个截面的族同构
	\[
		C\simeq\lambda^*\mathcal C,
	\]
	右边表示按 \(\lambda\) 拉回上述标准族。
\end{proposition}
\begin{proof}
	在一个平凡化邻域中，前三个截面给出三个随底点全纯变化且不相等的球面点。用前一命题的公式归一化；必要时先换一张球面坐标，使所用分母在当前底点非零，再缩小邻域。由公式可见归一化分式的系数随底点全纯变化，故给出局部族同构。第四截面在其中的坐标就是全纯函数 \(\lambda\)。交叠处的两个归一化都把前三截面送到 \(0,1,\infty\)，由唯一性逐纤维相等，因而作为全纯映射相等。它们遂粘合为全局同构及全局参数。由此得到交换图
	\[
		\begin{tikzcd}
			C \arrow[r] \arrow[d] & \mathcal C \arrow[d,"\pi"] \\
			T \arrow[r,"\lambda"'] & B.
		\end{tikzcd}
	\]
	上箭头在每条纤维上保持编号。唯一性同样归结为固定三个点的自同构恒等。
\end{proof}

\(B\) 不是单连通的：绕 \(0\)、\(1\) 的道路保留非平凡绕行信息。给球面四个点再指定拓扑标记，沿 \(B\) 中闭路移动第四点可能改变这个标记。前面定义的带标记空间记录这类信息；交比记录的是遗忘拓扑标记后的对象。上述族的构造和分类在整个 \(B\) 上成立，不需要选择一个交比的对数或一条回到基点的道路。

\subsection{标记碰撞、长颈与结点极限}
\label{g4:subsec:0605:03}
先令 \(\lambda=t\to0\)。在固定的 \(z\) 坐标中，第1和第4标记碰撞于零，极限不再有四个不同标记。为了保留两点的相对位置，引入放大坐标 \(\xi=z/t\)。这时两点始终为 \(\xi=0,1\)，而原先留在远处的部分从 \(\xi=\infty\) 接入。一个坐标适合远处，另一个适合碰撞的近处，它们需要在同一个曲面族中兼容。

\begin{proposition}
	\label{g4:prop:0605:node-family}
	取 \(|t|<\varepsilon<1/2\)。将 \(\mathbb P^1_z\times\{|t|<\varepsilon\}\) 的点 \((0,0)\) 替换为记录方向的射影直线，局部用
	\[
		zV=tU,\qquad [U:V]\in\mathbb P^1
	\]
	表示。所得复曲面光滑，投影到 \(t\) 为固有映射。\(t\ne0\) 的纤维仍是原四点球面；\(t=0\) 的纤维是两个球面在一个结点相交，每个球面恰含两个标记及该结点。结点附近的族方程为 \(xy=t\)。
\end{proposition}
\begin{proof}
	在 \(V\ne0\) 的坐标中令 \(\xi=U/V\)，方程为 \(z=t\xi\)，总空间以 \((t,\xi)\) 为光滑坐标。在 \(U\ne0\) 的坐标中令 \(\eta=V/U\)，方程为 \(t=z\eta\)，总空间以 \((z,\eta)\) 为光滑坐标。交叠处 \(\eta=1/\xi\)、\(z=t\xi\)，它们与未修改区域中的原坐标全纯相容。投影到原乘积是闭的射影方向条件，且方向因子紧，故为固有映射；再复合原来的固有投影即得所需固有性。对 \((z,t)\ne(0,0)\)，方向唯一，故纤维在 \(t\ne0\) 时不变。

	在第二张图中，中心纤维为 \(z\eta=0\)。\(z=0\) 是新增的方向球面，\(\eta=0\) 是原中心球面去掉零点后的闭包，两者在 \((z,\eta)=(0,0)\) 相交。第一张图中的 \(t=0\) 补全新增球面，其 \(\xi=\infty\) 就是这个交点。因此结点方程为 \(xy=t\)，取 \(x=z\)、\(y=\eta\)。

	第1、第4截面在第一张图中分别为 \(\xi=0\)、\(\xi=1\)，中心点各不相同，也不在结点。第2、第3截面仍位于原球面上的 \(z=1,\infty\)。于是各分支均有三个不同的特殊点，即两个标记和一个结点。任何保持所有标记的自同构必须保持两个分支，并在每个分支固定三个点，故为恒等。
\end{proof}

在结点图中固定 \(0<c<1\)，当 \(0<|t|<c^2\) 时，\(|x|<c\)、\(|y|<c\)、\(xy=t\) 的部分等价于
\[
	|t|/c<|x|<c.
\]
由\eqref{g4:eq:0505:plumbing-modulus}，其模量为 \(\log(c^2/|t|)/(2\pi)\)。所以碰撞不仅产生一个形式上的奇点，也产生模量趋于无穷的实际环形颈部。若 \(t=|t|e^{\mathrm{i}\theta}\)，粘合关系给出 \(\arg y=\theta-\arg x\)；辐角控制两端的相对转动，而模控制颈部的长度尺度。

\subsection{两两分配标记的三种稳定边界类型}
\label{g4:subsec:0605:04}
对于亏格零的结点曲线，稳定性的要求是：每个球面分支上的特殊点总数至少为三；特殊点包括标记点以及结点在该分支上的原像。这个条件正好排除固定一两个点仍可连续变化的球面自同构。

\ProseParagraphBreak

结点曲线的对偶图以归一化分支为顶点、结点为边。其算术亏格在这个模型中为各分支亏格之和加上对偶图的圈数。这个数也等于把各结点换成小颈部后光滑曲面的亏格：每挖去一个结点原像的小盘，Euler 示性数减少1；沿圆周接入圆柱不再改变它。对有 \(v\) 个分支、\(e\) 个结点的连通曲线，所得等式为 \(g=\sum g_i+e-v+1\)。所以算术亏格零恰要求所有分支为球面、对偶图为树。

\begin{proposition}
	\label{g4:prop:0605:three-boundaries}
	连通、算术亏格零、带四个有序光滑标记的稳定结点曲线，若非光滑，则恰由两个球面组成，各带两个标记，并在一个结点相接。其同构类型恰有三种，分别由无序二二分割
	\[
		\{1,4\}\mid\{2,3\},\qquad
		\{2,4\}\mid\{1,3\},\qquad
		\{3,4\}\mid\{1,2\}
	\]
	确定，它们分别对应 \(\lambda=0,1,\infty\)。
\end{proposition}
\begin{proof}
	亏格零意味着归一化的分支都是球面，连接各分支的对偶图是树。设有 \(v\) 个分支，则有 \(v-1\) 个结点。特殊点在各分支上的总数为 \(4+2(v-1)=2v+2\)。每个分支至少有三个，故 \(3v\le2v+2\)，即 \(v\le2\)。非光滑时 \(v=2\)，两边各有一个结点原像，稳定性遂迫使各自恰有两个标记。每个分支的三个特殊点都可唯一归一化为 \(0,1,\infty\)，所以一旦指定二二分割，曲线的同构类及其保持标记的同构就都确定。

	\(\lambda=0\) 的分割已由前一命题得到。\(\lambda\to1\) 时取 \(t=\lambda-1\)、\(x=z-1\)，碰撞的是第2和第4标记；重复同一坐标构造得到第二种。\(\lambda\to\infty\) 时取 \(t=1/\lambda\)、\(x=1/z\)，碰撞的是第3和第4标记，得到第三种。三个局部构造在非零参数处都恢复原族，因此可与 \(B\) 上的标准族粘合，给出参数球面上的稳定解析族。
\end{proof}

% !TeX root = ../Book4.tex
\begin{figure}[htbp]
	\centering
	\begin{tikzpicture}[scale=0.9, font=\small]
		\foreach \x/\a/\b/\c/\d/\v in {0/1/4/2/3/0,4.5/2/4/1/3/1,9/3/4/1/2/\infty}{
			\begin{scope}[xshift=\x cm]
				\draw[thick,GeometryBlue] (0,0)--(1.4,0);
				\draw (0,0)--(-0.6,0.7) node[above]{\(\a\)};
				\draw (0,0)--(-0.6,-0.7) node[below]{\(\b\)};
				\draw (1.4,0)--(2,0.7) node[above]{\(\c\)};
				\draw (1.4,0)--(2,-0.7) node[below]{\(\d\)};
				\fill[GeometryBlue] (0,0) circle (2pt) (1.4,0) circle (2pt);
				\node[below] at (0.7,-1.2){\(\lambda=\v\)};
			\end{scope}
		}
	\end{tikzpicture}
	\caption[有序四点球面的三个稳定边界]{三个边界纤维的对偶图。每个球面分支有两个标记和一个结点原像，三种分割由标记的编号区分。}
	\label{g4:fig:0605:boundary}
\end{figure}


图中的顶点代表球面分支，连接边代表结点，短枝的编号代表标记。它显示的是对偶图，不是将复曲线嵌入平面的图像。每个顶点恰连着三条枝或边，与稳定性条件相符。

\subsection{有序标记与置换商、自同构的区别}
\label{g4:subsec:0605:05}
保留四点的次序时，每个光滑对象的自同构群都平凡；若只记四点集合，自同构可以置换这些点。忘掉编号既改变参数的等价关系，也改变对象的对称群。

\begin{proposition}
	\label{g4:prop:0605:permutations}
	重排四个标记产生的交比恰属于
	\[
		\lambda,\quad 1-\lambda,\quad\lambda^{-1},\quad
		(1-\lambda)^{-1},\quad\frac{\lambda}{\lambda-1},\quad
		\frac{\lambda-1}{\lambda}.
	\]
	由 \(S_4\) 到这六个分式变换组成的群的同态，核为四元群
	\(V_4=\{1,(12)(34),(13)(24),(14)(23)\}\)。无序四点球面的自同构群通常有4个元素；当交比属于 \(\{-1,1/2,2\}\) 时有8个元素，当 \(\lambda^2-\lambda+1=0\) 时有12个元素。
\end{proposition}
\begin{proof}
	交换第1与第2点，把标准配置重新归一化，得到 \(1-\lambda\)；交换第2与第3点得到 \(\lambda/(\lambda-1)\)。代入公式还得到逆数变换。这些变换相互复合只有所列六个有理函数，并在三个缺失值上实现全部置换，所以它们组成 \(S_3\)。逐项代入三个双换位，交比均不变。群同态满射且核有 \(24/6=4\) 个元素，故核正是 \(V_4\)。

	固定一个配置后，能由 Möbius 变换实现的点置换，恰是保持该配置交比的置换；存在性来自归一化，唯一性来自前三点。其数目等于 \(4\) 乘以 \(\lambda\) 在上述 \(S_3\) 作用下的稳定子阶数。三个非平凡对合的固定点在 \(B\) 内分别为 \(1/2,-1,2\)，两个三阶变换的固定点满足 \(\lambda^2-\lambda+1=0\)。两类集合不交，故稳定子阶数分别为 \(2,3\)，其他点为 \(1\)，得到所述计数。
\end{proof}

例如交比为 \(-1\) 的配置具有额外对称，因而无序参数的商映射在那里出现二重分歧。它仍是光滑四点球面，绝不是结点边界。有限对称造成的商分歧，与 \(\lambda\to0,1,\infty\) 时产生新分支的退化，是两种不同的几何现象。

\subsection{与 V.34.6 的代数通用族及完整紧化构造比较}
\label{g4:subsec:0605:06}
上述计算得到一个完整的一维解析模型：\(B\) 参数化光滑有序四点球面；补上三个参数值，得到三种稳定结点纤维。局部 \(xy=t\) 将参数的消失、环域模量发散与核心圈收缩联系起来。保持标记的自同构在光滑纤维及三个边界纤维上均平凡，因此这些具体对象没有因对称性产生的识别歧义。

\ProseParagraphBreak

第五卷第34.6节将把同一模型放在代数基底上：从 \(\mathbb P^1_z\times\mathbb P^1_\lambda\) 的四个截面出发，在三个碰撞中心作吹起，并验证平坦性、基变换与通用性。本节的方向坐标已经给出了复数域上相应的局部方程。代数的通用性还要求对非约化测试基底构造拉回及下降，因此除了几何点分类，还需检查无穷小底空间上的族；这些是第五卷的具体证明任务。
% 纲要条目排为编号小节标题，并进入目录与书签；说明段仍保留于纲要注释中。
